\chapter{Problem Formulation and Experimental System}
\label{ch:system}

\chapterhandoff{Chapter~\ref{ch:bg} showed that no existing line of work
chooses the page representation for an otherwise unchanged agent, using only
what can be known before the expensive option is paid for. This chapter
defines the system in which such a choice can be measured. Its purpose
is to make one thing auditable: which differences observed later are
attributable to representation, and which are not.}

\section{The comparison boundary}
\label{sec:boundary}

Every empirical claim in this dissertation rests on a single structural
commitment, and it is worth stating before any implementation detail. Across
all conditions we hold constant: the task and the state of the website it
starts from, the underlying model and the settings that govern how it produces
text, the set of browser actions available, the limit on how many steps are
allowed, the browser and the size of its window, the benchmark's automatic
scoring, and the rules for what counts as a cost. We vary exactly one stage: the construction of the observation handed to the model at each step.

\begin{figure}[t]
  \centering
  \includegraphics[width=\textwidth]{fig_f3_comparison_boundary.pdf}
  \caption[The causal comparison boundary]{%
  \textbf{What is varied and what is held identical.} Five pipeline stages carry
  the episode from task to scored outcome. Four are grey and marked
  \emph{same}; only observation construction differs across conditions. The
  dashed boundary marks what is held identical, and cost is measured at the same boundary as outcome, so a mode that spends more to reach the same result shows up as more expensive rather than being absorbed. This is what licenses attributing
  later differences to representation; it is not an architecture diagram of the
  codebase.}
  \label{fig:f3}
\end{figure}

That is the whole design. Everything in the rest of this chapter is a matter of
making the commitment concrete and checkable.

\section{Task, episode, step}
\label{sec:task-def}

A \emph{task} is a benchmark-defined user objective on a specific website, with
an instruction written in ordinary language, a starting web address, sometimes
a reference image, and a scoring procedure that runs automatically. An \emph{episode} is one attempted execution of one task
from its initial state until success, an explicit stop action, an unrecoverable
error, or exhaustion of the step budget. A \emph{step} is one
observe--decide--act iteration.

At step $t$ the agent receives observation $o_t$, its interaction history $h_t$,
and emits an action $a_t$ from a fixed space: \texttt{click}, \texttt{type},
\texttt{select}, \texttt{hover}, \texttt{scroll}, \texttt{go\_back},
\texttt{go\_forward}, \texttt{goto}, \texttt{new\_tab}, \texttt{close\_tab},
\texttt{stop}. The action space does not vary with the observation mode. This
matters more than it may appear: if the image-carrying modes could emit
coordinate clicks while the text modes could not, mode differences would
partially be action-space differences, and no representation attribution would
survive. All modes ground actions through element identifiers.

\paragraph{Why a step budget of 30.} The budget must be large enough that
capable modes are not truncated and small enough that hopeless episodes do not
dominate the cost accounting. Thirty steps sits above the observed length
distribution of successful episodes on both VisualWebArena sites; measured mean
episode length across all conditions ranges from 13.7 to 28.4 steps
(Table~\ref{tab:cost-per-mode}), so the budget binds for the failing tail rather
than for the working modes. It is set identically for every condition, which is
the property the comparison actually needs.

\section{The controlled observation space}
\label{sec:obs-space}

An observation is assembled from three components, each of which we control
independently.

\paragraph{Text payload.} Either the accessibility-tree serialisation (AXTree),
which lists elements with roles, names, states and identifiers; or the
\sommarks{} listing, a flattened enumeration of interactable elements indexed in
the Set-of-Marks convention. These are two encodings of overlapping information,
and the difference is one of format rather than of quantity. Only step~0 is comparable: there the two runs are guaranteed to be looking at the same page, and trajectories diverge as soon as the modes choose different actions. Measured across every task at that step, the \sommarks{} payload is a median $1.008\times$ the
AXTree payload on classifieds ($n=224$ tasks, 97.3\% of pairs within $\pm5$\%)
and $1.004\times$ on reddit ($n=205$, 71.2\% within $\pm5$\%).

\begin{scopebox}
\textbf{Why only step~0.} Pooling the same ratio over all 2{,}506 mirrored
(task, step) pairs gives a median of 0.980 and a 5th--95th percentile range of
0.38--1.86. That dispersion is not payload variation: past step~0 the two runs
have taken different actions and are describing \emph{different pages}, so the
pooled figure compares things that are not comparable. The step-0 restriction is
what makes the ratio a measurement rather than an average over a confound.
\end{scopebox}

\paragraph{Prompt framing.} Either the structured-text instruction set, which
tells the model it is reading an accessibility tree and how to refer to elements;
or the Set-of-Marks instruction set, which tells the model that elements carry
visible numbered marks and how to refer to those.

\paragraph{Image.} Either no image, an annotated screenshot with numbered
overlays on candidate elements, or a raw un-annotated screenshot.

\subsection{The six observation modes}
\label{sec:modes}

Crossing text payload with prompt framing yields four image-free
configurations; adding the two image-bearing conditions gives six modes in
total (Table~\ref{tab:modes}).

\begin{table}[t]
\centering\small
\caption[The six observation modes]{\textbf{The six observation modes.} The four
image-free modes are a $2\times2$ of text payload $\times$ prompt framing; the
grouping in the last column is verified from the step records rather than asserted: all four text-side modes log
\texttt{image\_payload\_bytes}~$=0$.}
\label{tab:modes}
\begin{tabular}{@{}llllc@{}}
\toprule
Mode & Text payload & Prompt framing & Image & Side \\
\midrule
\mDOM{}   & AXTree   & structured-text & none             & text \\
\mPP{}    & AXTree   & Set-of-Marks    & none             & text \\
\mPT{}    & \sommarks{} & structured-text & none             & text \\
\mPS{}    & \sommarks{} & Set-of-Marks    & none             & text \\
\mSoM{}   & \sommarks{} & Set-of-Marks    & annotated       & combined \\
\mVis{}   & none     & vision-oriented & raw             & visual \\
\bottomrule
\end{tabular}
\end{table}

\paragraph{Why this factorisation exists.} The three-mode comparison that this study began with (structured text, Set-of-Marks, raw vision) cannot separate three changes that it makes simultaneously. Moving from \mDOM{} to
\mSoM{} changes the text encoding, the prompt, \emph{and} the image at once, so
any difference is uninterpretable as evidence about any one of them. The four
image-free modes exist to break that confound: \mPS{} is \mSoM{} with the image
removed and nothing else changed, \mPT{} isolates the text encoding, and \mPP{}
isolates the prompt. This is experimental refinement in response to an
identified ambiguity, not a proposal for a new agent design.

\paragraph{Naming.} We use \emph{phantom} for the three image-free modes other than \mDOM{}: \mPT{}, \mPP{} and \mPS{}. \mDOM{} is image-free too,
so ``text side'' and ``phantom'' are not synonyms, and we do not use them
interchangeably. \mPS{} is one such configuration; when we speak of the
\emph{phantom routing space} we mean the whole set of image-free
configurations that \mPT{}, \mPP{} and \mPS{} are drawn from.

\section{Routing formulation}
\label{sec:router-formulation}

Let $\mathcal{M}$ denote the set of observation modes. At step $t$ the router
observes a feature vector $z_t$ drawn only from information available
\emph{before} any mode-specific observation is constructed, and selects a mode:
\begin{equation}
  m_t = \pi_r(z_t), \qquad m_t \in \mathcal{M},
  \label{eq:router}
\end{equation}
after which the agent policy acts on the observation that mode produces:
\begin{equation}
  a_t = \pi_a\!\left(o_t^{(m_t)},\, h_t\right).
  \label{eq:agent}
\end{equation}
In words: the router picks how the page will be shown, and the agent then
picks what to do with what it sees. The formulation is written per step
because that is the general case, but every result in this dissertation uses
the task-level special case, in which the router chooses once and that mode is
used for the whole episode. The reason is supply, not principle: the benchmark
scores whole tasks, so only whole tasks carry an outcome a router could be
trained or judged on (Section~\ref{sec:regime-boundary}). The agent's decision rule $\pi_a$ is
identical in every condition; only the router $\pi_r$ changes. A \emph{fixed
policy} is the special case where the router always returns the same mode,
$\pi_r(z_t) \equiv m$.

\subsection{The oracle and its decomposition}

Let $s_i^{(m)} \in \{0,1\}$ be the success of task $i$ under mode $m$, and
$c_i^{(m)}$ its measured cost. The retrospective oracle is
\begin{equation}
  \mathrm{SR}_{\text{oracle}} = \frac{1}{n}\sum_{i=1}^{n} \max_{m\in\mathcal{M}} s_i^{(m)},
  \label{eq:oracle}
\end{equation}
which is an upper bound rather than an achievable result: for each task it
takes the best outcome across all six modes, including modes a real system
would not have run. Throughout this dissertation the \emph{ceiling} means this
quantity's gain over the best single mode, $\mathrm{SR}_{\text{oracle}} -
\max_m \mathrm{SR}^{(m)}$, always reported against a stated comparator.

The oracle splits into two halves that differ sharply in how many labelled
examples they require to learn, and Chapter~\ref{ch:pred} depends on keeping them apart:
\begin{itemize}[leftmargin=1.5em]
\item \textbf{Triage.} For each task, decide whether \emph{any} mode will
      solve it; if not, spend the cheapest. The label is binary and is defined
      on every task, so its supply is never scarce.
\item \textbf{Which-mode.} Among tasks that some mode solves, decide which.
      The label exists only where a success occurred, so its supply is bounded
      by the success rate itself.
\end{itemize}

\subsection{Drop-one contribution}

The unique contribution of mode $m$ to the oracle portfolio is what removing it
costs:
\begin{equation}
  \Delta_m = \frac{1}{n}\sum_{i=1}^{n}
    \left[\max_{m'\in\mathcal{M}} s_i^{(m')} - \max_{m'\in\mathcal{M}\setminus\{m\}} s_i^{(m')}\right].
  \label{eq:dropone}
\end{equation}
$\Delta_m>0$ means $m$ solves at least one task nothing else solves.
Equation~\ref{eq:dropone} also describes, in the same form, what one extra
\emph{repeat} of a mode already in the set would contribute. That is exactly why Chapter~\ref{ch:repr}
can weigh a new representation against a repeat of an old one while giving
both sides the same number of runs.

\subsection{Net saving}

For a cost variable $C$ stated at the point of use,
\begin{equation}
  \mathrm{NetSaving} = C_{\text{fixed comparator}} -
    \bigl(C_{\text{routed execution}} + C_{\text{router overhead}}\bigr),
  \label{eq:netsaving}
\end{equation}
where the overhead term covers feature extraction, the routing model's own
inference, any additional observation construction, and retries attributable to
routing. Reporting the first two terms without the third is not an evaluation of
a policy.

\section{Benchmarks, sites and scope}
\label{sec:benchmarks}

\paragraph{Primary setting.} VisualWebArena \citep{koh2024visualwebarena}
provides self-hosted, deterministic-evaluator web environments with
visually-grounded tasks. We use two sites, \emph{classifieds} and
\emph{reddit}. Both are self-hosted from the published images, so the
environment does not drift under us in the way a live website would.

\paragraph{External validation.} WebArena \citep{zhou2024webarena} reddit is
used to test whether the findings survive on a benchmark whose tasks are
specified differently, for two of the models. Its value here is
specific and measurable rather than rhetorical: WebArena task
configurations \emph{omit the reference-image field entirely}, all 480 of them across its three sites, whereas 29--40\% of VisualWebArena tasks supply
one (Figure~\ref{fig:f4}, panel B). The distinction between an absent field and
an explicit zero is worth keeping: what we observe is a schema difference, from
which no reference image follows, and not an annotation asserting that none is
needed.

\begin{scopebox}
\textbf{A caveat we impose on ourselves.} It is tempting to read ``no reference
images'' as ``the visual demand that motivates Set-of-Marks is absent''. That conflates \emph{visual matching} (comparing the page against a supplied target image, which is genuinely gone on WebArena) with \emph{visual grounding} (locating and acting on elements of the current screen, which is what the marks actually serve and which is still required). We therefore claim
only a difference in task specification, not a difference in modality
requirement. Establishing the stronger claim would need evidence that a
text-only mode closes the gap to \mSoM{} on WebArena but not on VisualWebArena,
which we do not have.
\end{scopebox}

\paragraph{Excluded sites.} VisualWebArena \emph{shopping} was run in an early
phase but is not part of the primary analysis; its conditions predate several
pipeline corrections and mixing them with post-correction data would compare
different measurement regimes. Chapter~\ref{ch:disc} states what its absence
costs the external-validity argument.

\subsection{From corpus to scored set}
\label{sec:scored-set}

Exclusions apply in two layers at different times, and the distinction is operational rather than cosmetic.

\emph{N/A tasks} are removed at task-load time: these are string-match tasks
whose reference answer is the literal ``N/A''. Our agent prompt provides no
``this is impossible'' exit, so they are unpassable by construction, and the
evaluator cannot distinguish a reasoned N/A judgement from an early exit. They
are removed before the runner ever sees them.

\emph{Protocol exclusions} are removed at analysis time only. The runner keeps
collecting them deliberately, so the exclusion stays auditable and reversible.

\begin{table}[t]
\centering\small
\caption[Corpus, run set and scored set]{\textbf{Two counts are both correct,
and this dissertation uses both.} The \emph{run set} is what a completed run
contains; the \emph{scored set} is the denominator of every reported success
rate. This is the source of the twin numbers that appear in our logs (reddit as
both 205 and 203).}
\label{tab:scored-set}
\begin{tabular}{@{}lrrrrr@{}}
\toprule
Corpus & size & $-$N/A & run set & $-$protocol & \textbf{scored set} \\
\midrule
VWA classifieds & 234 & 10 & 224 & 0 & \textbf{224} \\
VWA reddit      & 210 &  5 & 205 & 2 & \textbf{203} \\
WA reddit       & 106 &  2 & 104 & 0 & \textbf{104} \\
\bottomrule
\end{tabular}
\end{table}

The denominator of every success rate in this dissertation is the scored set, and within a cell all conditions are evaluated on the byte-identical task-identifier set, verified by SHA-256 of the sorted identifier list. The analysis pipeline refuses any condition whose set fails that check, rather than silently truncating it.

\subsection{What the corpora actually contain}

\begin{figure}[t]
  \centering
  \includegraphics[width=\textwidth]{fig_f4_corpus_eda.pdf}
  \caption[Corpus characterisation across the six task sets]{%
  \textbf{The corpora these conclusions rest on.} (A) The run-set/scored-set
  distinction of Table~\ref{tab:scored-set}. (B) What each corpus annotates: WebArena carries neither reference images nor difficulty labels, which
  Chapter~\ref{ch:pred} shows is not incidental. (C) Evaluator types.
  (D) Instruction length. Two typos in the distributed annotations
  (\texttt{mediun}, \texttt{hrad}) are reported rather than silently repaired,
  because folding a typo into the nearest label would make a
  difficulty-stratified analysis run on a different partition than the one it
  documents.}
  \label{fig:f4}
\end{figure}

Two properties of the corpora bear on later interpretation. First, task
diversity is well below task count: VisualWebArena classifieds' 234 tasks are
generated from 75 distinct templates, and WebArena reddit's 106 from 21. The
template count is the honest denominator for ``how many different things are we
asking''. Second, the annotation asymmetry in Figure~\ref{fig:f4}(B) reappears
in Chapter~\ref{ch:pred} as the single most consequential fact about the
feature set.

\section{Models and inference configuration}
\label{sec:models}

Three backbones span two model families and two deployment classes
(Table~\ref{tab:backbones}).

\begin{table}[t]
\centering\small
\caption[Backbone models]{\textbf{Three backbones, two families, two deployment
classes.} B1 and B2 are matched in scale and differ in family, which is what
makes B2 a cross-family control rather than a third data point.}
\label{tab:backbones}
\begin{tabular}{@{}llll@{}}
\toprule
Label & Model & Serving & Role \\
\midrule
B0 & \texttt{Qwen3-VL-235B-A22B} & metered API proxy & capability upper end \\
B1 & \texttt{Qwen3-VL-4B}        & local, bf16       & small-model reference \\
B2 & \texttt{google/gemma-3-4b-it} & local, bf16     & cross-family control at matched scale \\
\bottomrule
\end{tabular}
\end{table}

\paragraph{Decoding.} All three run greedy (\texttt{temperature}~$=0$,
\texttt{do\_sample}~$=$~false) with \texttt{max\_new\_tokens}~$=4096$, seed 42.
Greedy decoding is chosen so that observed run-to-run variation is attributable to the environment rather than to sampling, a property we exploit in
Section~\ref{sec:noise-floor} and, importantly, one we \emph{test} rather than
assume: step-level determinism does not imply episode-level determinism, because an episode also depends on site state and session lifetime.

\paragraph{Why B2 is a control and not a third result.} An earlier framing of
this study treated the interaction between backbone capability and
representation as a third contribution in its own right, on the strength of a
large B0-versus-B1 gap. That gap turned out to be a capability difference that
appears under \emph{every} representation rather than an interaction with
representation, so the framing was dropped and the small backbones were
re-scoped as controls. B2 is included because it differs from B1 in family while matching it in scale; it is not a third point on a capability curve.

\paragraph{Why a 4B model at all.} B1 and B2 succeed on 0.45--14.29\% of tasks
depending on site and mode. A reader may reasonably ask what a routing study
learns from backbones this weak. The answer is that the base success rate is not a nuisance parameter here. It is the independent variable that
Chapter~\ref{ch:pred} identifies as the binding constraint. Having cells that
span 0.45\% to 35.58\% single-mode success is what allows the label-supply
mechanism to be observed as a gradient rather than asserted at one operating
point.

\paragraph{The proxy caveat, stated where it applies.} B0 is served through a
metered API proxy whose protocol is a hybrid: an Anthropic-style endpoint and
header with an OpenAI-style tool schema and top-level tool-call and log-prob
fields. This is an implementation fact with two consequences we carry forward.
First, B0 conditions are replicable from our archived logs but not
independently reproducible without equivalent proxy access. Second, the
service reports how confident the model was only for the two most likely next
words rather than for all of them. Four of our six confidence measures can
therefore be computed for B0 and two cannot. The inputs used in
Chapter~\ref{ch:pred} do not depend on the two that are missing.

\section{Metrics and the accounting boundary}
\label{sec:metrics}

\paragraph{Primary outcome.} Task success as returned by the benchmark
evaluator, aggregated as task success rate (TSR) over the scored set. The
evaluator is unchanged from the benchmark distribution.

\paragraph{Primary cost.} \texttt{total\_billed\_cost\_usd}: the per-episode sum
over steps of input and output tokens, priced at that backbone's per-1k
schedule. For B0 the schedule is the proxy's actual billing rate, so the figure
is money charged. For B1 and B2 the schedule is a device-amortisation rate, so
the figure is a \emph{token-priced cost estimate} rather than an invoice.
These are two different \emph{deployment classes} --- money actually charged,
and an estimate of what the hardware time was worth --- and we compare costs
only within one of them, never across. This single quantity is used in every
policy comparison in Chapters~\ref{ch:repr}--\ref{ch:e2e}.

\begin{scopebox}
\textbf{Every cost number in this dissertation is governed by two rules.} \emph{(i)}
Comparisons are made \emph{within} a cell only. Across the API/local boundary
the schedules are different objects, and the ratio between them (82--107$\times$
against measured electricity) is a statement about deployment mode, not model
efficiency. \emph{(ii)} Token-priced cost and measured device energy are
\emph{different quantities} and are never substituted for one another. The
policy tables use the former; Section~\ref{sec:sustainability} and only that
section uses the latter.
\end{scopebox}

\begin{table}[t]
\centering\small
\caption[Cost per mode, canonical estimand]{\textbf{Mean per-episode cost
(\texttt{total\_billed\_cost\_usd}) by mode and cell.} Cheapest mode per column
in bold. What matters downstream is this: the cheapest mode is not the same
everywhere (\mDOM{} on the WebArena forum with the 235B model, \mVis{}
elsewhere), and the spread between cheapest and dearest ranges from 11.6\% to
77.8\%, so ``the expensive mode'' means something different in each cell. Cell
codes name a site and a model: \cell{cls} and \cell{red} are the
classified-ads and forum sites of VisualWebArena, \cell{wa\_red} the WebArena
forum; \cell{B0} is the 235B model, \cell{B1} and \cell{B2} the two 4B models,
one from each design family.}
\label{tab:cost-per-mode}
\footnotesize\setlength{\tabcolsep}{4.2pt}
\begin{tabular}{@{}lrrrrrrrr@{}}
\toprule
Mode & \cell{cls$\cdot$B0} & \cell{red$\cdot$B0} & \cell{cls$\cdot$B1}
     & \cell{red$\cdot$B1} & \cell{cls$\cdot$B2} & \cell{red$\cdot$B2}
     & \cell{wa\_red$\cdot$B0} & \cell{wa\_red$\cdot$B1} \\
\midrule
\mDOM{}  & 0.0696 & 0.1015 & 0.0595 & 0.0733 & 0.0768 & 0.0948 & \textbf{0.0753} & 0.0658 \\
\mPP{}   & 0.0685 & 0.1016 & 0.0630 & 0.0766 & 0.0845 & 0.0994 & 0.0775 & 0.0739 \\
\mPT{}   & 0.0692 & 0.1058 & 0.0588 & 0.0695 & 0.0732 & 0.0885 & 0.0848 & 0.0615 \\
\mPS{}   & 0.0721 & 0.1081 & 0.0597 & 0.0748 & 0.0846 & 0.0945 & 0.0850 & 0.0666 \\
\mSoM{}  & 0.0724 & 0.1105 & 0.0603 & 0.0800 & 0.0908 & 0.1116 & 0.0911 & 0.0794 \\
\mVis{}  & \textbf{0.0648} & \textbf{0.0981} & \textbf{0.0432} & \textbf{0.0524}
         & \textbf{0.0707} & \textbf{0.0683} & 0.0864 & \textbf{0.0447} \\
\midrule
spread & $+11.6$\% & $+12.6$\% & $+46.1$\% & $+52.7$\%
  & $+28.5$\% & $+63.3$\% & $+21.0$\% & $+77.8$\% \\
\bottomrule
\end{tabular}
\end{table}

\paragraph{Secondary measures.} Step count (mean 13.7--28.4 across all
conditions); wall-clock latency, decomposed into model inference and environment
interaction; device energy from NVML sampling at 0.5\,s for the locally served
backbones; and a page-change diagnostic. We note explicitly that
\texttt{page\_changed}~$=$~false means \emph{no detected page change} under our
detector, not \emph{action failed}; the two are routinely conflated and are kept
separate here.

\paragraph{Measured energy, kept separate.} For B1 and B2 we additionally
measure device energy directly, which yields a quantity the token schedule
cannot: on classifieds, 0.00482--0.00777\,kWh per episode, and on reddit
0.00830--0.01236\,kWh. Notably the energy ordering does not always match the token-priced ordering: on \cell{cls$\cdot$B2}, \mVis{} is the cheapest mode by
tokens and the \emph{most} energy-intensive by measurement, because a weaker
backbone spends more wall-clock steps in it. That divergence is exactly why the
two are not treated as proxies for one another.

\subsection{What our energy measurement is and is not}
\label{sec:sustainability}

We measure GPU device power via NVML at 0.5\,s intervals during locally-served
inference and integrate it over the episode. That yields \emph{GPU-device
operational energy}. Multiplying by a stated grid carbon intensity (220\,g
CO\textsubscript{2}e/kWh, UK average) yields a \emph{GPU-device operational
emissions estimate}.

It does not yield a carbon footprint. Excluded are: host CPU and memory, cooling
and facility overhead (no PUE is applied), network, the browser environment's own
compute, embodied emissions of the hardware, and (for B0 entirely) the remote data-centre's consumption, which we cannot observe. Following
\citet{gupta2022chasing} and \citet{luccioni2022bloom}, the boundary chosen
materially determines the number, so we name ours and decline the broader label.
For the same reason, no token-count reduction anywhere in this dissertation is
converted into an energy or emissions reduction: no hardware-independent
conversion exists \citep{fernandez2025energy}.

\section{Statistical design}
\label{sec:stats}

\paragraph{Pairing.} Every mode within a cell runs the same task identifiers, so
all within-cell comparisons are paired at the task level, and all oracle and
drop-one quantities are computed on the complete-case intersection.

\paragraph{Stratification.} The unit of stratification is the \emph{cell} =
(site, model). The unit of operational launch is the \emph{condition} = (site,
model, mode). Six modes $\times$ 6 VisualWebArena cells $=36$ primary
conditions, plus 12 WebArena conditions across 2 cells
(Figure~\ref{fig:f5}). These two words are not interchangeable and we do not mix
them.

\paragraph{Uncertainty.} Drop-one contributions carry bootstrap 95\%
intervals over tasks. Because every mode attempts the same tasks, differences
between two modes are tested by counting the tasks on which exactly one of
them succeeded (McNemar's test). Where a family of tests is reported over cells, we apply Holm's
step-down correction \citep{holm1979} with the family size stated explicitly at
the point of use.

\paragraph{The floor that governs everything.} No effect size in this
dissertation is interpreted without reference to a measured same-condition
noise floor, obtained by rerunning identical conditions and comparing the
resulting solve sets. Section~\ref{sec:noise-floor} derives it. This is the
single methodological decision that most changes the conclusions, and it changes
them in the direction of weakening our own positive claims.

\begin{figure}[t]
  \centering
  \includegraphics[width=\textwidth]{fig_f5_design_matrix.pdf}
  \caption[The experimental design matrix]{%
  \textbf{What was run.} Rows are observation modes, columns are (site, model)
  cells; the two rightmost columns are the WebArena external-validation cells.
  Each filled square is one operational condition; each column is one
  statistical stratification unit. The grid has no holes, which is what makes
  every within-cell comparison paired. Completeness against the expected task-ID
  set is enforced upstream by the aggregators
  (Section~\ref{sec:repro}), not re-verified by the figure script.}
  \label{fig:f5}
\end{figure}

\section{Reproducibility}
\label{sec:repro}

Each step writes a structured JSONL record carrying the observation
configuration, token counts, latencies, action, execution result and
page-change diagnostic; each episode writes a summary with the evaluator's
verdict. Analysis reads these records through a single deduplicating loader, so
a restarted run cannot double-count. The runner and the analysis pipeline are
separate processes with no shared state, and the figure scripts read only
committed analysis artefacts.

The pipeline has two habits that exist because their absence caused errors, and they are worth naming since they affect what a reader can check. Conditions are refused
by the aggregators unless their task-identifier SHA matches the cell's expected
set, which converts a silent partial run into a loud failure. And every analysis artefact records a provenance header naming the producing script, the
capture timestamp, and whether the file was approved for use in the main
analysis or produced only for exploration, so a number lifted into prose can
be traced back to the run that produced it.
Appendix~\ref{app:repro} lists the configuration manifest, model revision pins
and artefact inventory.

\section{Keeping the study honest}
\label{sec:honesty}

A negative result is worth exactly as much as the discipline that produced it. A
reader who cannot see how we guarded against fooling ourselves has no reason to
prefer our ``routing does not pay here'' to the more ordinary explanation that
we measured badly. This section states the four rules that did that work, and it
is deliberately explicit about which of them we fixed in advance and which we adopted only after something went wrong, because that split is itself part of the answer.

\paragraph{Fixed in advance, and externally witnessed.} The hypotheses, the
primary gate, the analysis plan and the exclusion policy were written before the
paper-grade data existed, and locked on 2026-05-18 by a repository tag and an
OSF registration DOI issued while the data-collection status still read
\emph{pending}. The protocol did not survive contact with the experiment unchanged: it was revised ten times. Each revision is a separate dated tag rather than a silent edit, so the difference between what we planned and
what we did is a matter of public record instead of memory. That is the property
that matters; an unamended preregistration on a four-month project would be less
credible, not more.

\paragraph{Fixed in advance: refuse rather than truncate.} Aggregators compare
the SHA-256 of a condition's sorted scored task-identifier list against the
cell's expected set and refuse anything that does not match, rather than
computing a rate over whatever episodes happen to be present. A partial run therefore stops with an error instead of quietly shrinking a denominator.

\paragraph{Adopted mid-project, after it caught us out: judge every effect
against a measured repetition floor.} Nothing in the original plan required
comparing an effect size to what rerunning an unchanged condition would produce.
We introduced the requirement after replicating conditions revealed 10--14\%
task-level disagreement between identical runs (Section~\ref{sec:noise-floor}),
and it immediately cost us a result: the complementarity structure that an
earlier version of this work advanced as a positive finding does not clear it
(Section~\ref{sec:dropone}). The rule has since taken the ceiling claim down
with it, from a statement about how high the ceiling is to one about what it
costs (Section~\ref{sec:six-vs-six}).

\paragraph{Adopted mid-project: adversarial audit by readers with no stake.}
Before each milestone the draft and the code behind it were read cold by three
AI systems built by three different developers, each writing its own
claim-by-claim attack without seeing the others' output. Their disagreements were the useful part, since the overlap was mostly things we already knew. Every
finding was then checked against the artefacts by hand before being accepted.
That was not ceremony. Of the attacks that reached a fix, one had its own
arithmetic wrong --- it reported a significance threshold of 4.91\pp{} where
the exact calculation gives 4.02\pp{} --- and two were factually wrong about
our pipeline. One of those claimed we had dropped a feature after noticing it
performed too well, when in fact the benchmark it had to transfer to does not
carry that feature at all.

\paragraph{What this bought, and what it does not guarantee.} The corrections
these rules produced are not confined to a footnote; they are distributed through
the chapters that follow, and are marked where they occur. They include the cost
range in Section~\ref{sec:ceiling}, the significance threshold in
Section~\ref{sec:floor-spread}, the complementarity claim in
Section~\ref{sec:dropone}, the determinism assumption in
Section~\ref{sec:noise-floor}, the cross-validation protocol in
Section~\ref{sec:protocol}, the ``untrainable'' phrasing in
Section~\ref{sec:which-half}, and the labelling rule in
Section~\ref{sec:mechanism}. In every case the corrected version is weaker than
the one it replaced.

None of this makes the result true. A discipline that catches errors is evidence
about the process, not about the world, and the errors it did not catch are by
construction invisible to it. What it does establish is that the claims which survived were not the ones we wanted, which is the only form of assurance a
negative result can offer.

\chapterhandoff{This chapter fixes the terms of the comparison: the task, the
agent, the actions it may take and the way costs are counted are held
identical, and only the way the page is presented changes. That boundary makes representation differences interpretable, but it does not make them \emph{interesting}. If one mode dominated everywhere,
selection would have nothing to select. Chapter~\ref{ch:repr} therefore asks the
prerequisite empirical question: do these modes succeed and fail on different
tasks, and if so, how much is that worth?}
